\documentclass[11pt]{article}
\usepackage[a4paper,margin=18mm]{geometry}
\usepackage{fontspec}
\setmainfont{Latin Modern Roman}
\usepackage{amsmath,amssymb,amsthm,amscd,graphicx,color}
\usepackage{xurl}
\usepackage[unicode,hidelinks]{hyperref}
\allowdisplaybreaks
\setlength\emergencystretch{3em}
\usepackage{amscd}
\usepackage[all]{xy}

\newtheorem{thm}{Theorem}[section]
\newtheorem{prop}[thm]{Proposition}
\newtheorem{lem}[thm]{Lemma}
\newtheorem{cor}[thm]{Corollary}
\newtheorem{claim}[thm]{Claim}

\theoremstyle{definition}
\newtheorem{dfn}[thm]{Definition}
\newtheorem{rem}[thm]{Remark}

\newcommand{\C}{\mathbb{C}}
\newcommand{\R}{\mathbb{R}}
\newcommand{\Z}{\mathbb{Z}}

\newcommand{\pt}{\mathrm{pt}}



\makeatletter
\def\tsep#1{%
  {\@tempdima=\ht\strutbox\advance\@tempdima #1
   \vrule height \@tempdima depth 0pt width 0pt\nobreak\hspace{0pt}%
  }%
}
\def\bsep#1{%
  {\@tempdima=\dp\strutbox \advance\@tempdima #1
   \nobreak\hspace{0pt}\vrule height 0pt depth \@tempdima width 0pt
  }%
}
\makeatother
\makeatletter\newlabel{thm:main}{{1.1}{}{Original source locator}{}{}}\makeatother
\begin{document}
\begin{center}\Large Integral Borel cohomology through degree three and its second filtration for the order-twelve dihedral action on the hexagonal torus\end{center}
\noindent\textbf{Stable ID:} 1882\par
\noindent\textbf{Authors:} Kiyonori Gomi\par
\noindent\textbf{Original work:} Twists on the Torus Equivariant under the 2-Dimensional Crystallographic Point Groups\par
\noindent\textbf{Version:} Official SIGMA 13 (2017) article 014, DOI 10.3842/SIGMA.2017.014; journal PDF and native TeX acquired 2026-10-01, exact bytes pinned by SHA256\par
\noindent\textbf{Selected task scope:} Only the hexagonal lattice Pi=\allowbreak{}Z(1,0)+\allowbreak{}Z(1/2,sqrt(3)/2), acted on by C=\allowbreak{}rotation through pi/3 and sigma=\allowbreak{}reflection in the first coordinate axis, with D6 of order12. Theorems4.11 and4.14 together give H0=\allowbreak{}Z,H1=\allowbreak{}0,H2=\allowbreak{}(Z/2)\textasciicircum{}2,H3=\allowbreak{}(Z/2)\textasciicircum{}2 and F2H3=\allowbreak{}Z/2. No other wallpaper action, twisted coefficients or global table is selected.\par
\noindent\textbf{Difficulty:} H2: Computing this Borel equivariant cohomology requires a stabilizer-sensitive cell decomposition and two Mayer-Vietoris sequences with explicit character restriction maps. Recovering the Leray-Serre filtration additionally requires a nontrivial integral module extension and the multiplication-by-three map in group cohomology. These constructions and their compatibility exceed a routine application of the spectral sequence or ordinary qualifying topology.\par
\noindent\textbf{Editorial boundary note (not author text):} One unified p6m/D6 Borel-cohomology and Leray-Serre-filtration computation is selected, containing both complete original Theorems 4.11 and 4.14. The exact hexagonal lattice/action, full native editable CW diagrams, stabilizer-sensitive two Mayer–Vietoris character restriction calculations, free-cell excision and integral-module extension whose composition is multiplication by three are included. The full fixed-point/sign-coefficient and low-degree twisted-point support is bundled as uncounted core; the source remaining low-degree exact-sequence calculations are expressly routine and retained as stated. Known finite-group integral cohomology is cited background. No other wallpaper/twisted-coefficient variant or stable-splitting assertion is counted. The original D4/D2 local character-group naming slip is unchanged with its locally explicit subgroup/maps.\par
\noindent\textbf{Source:} \url{https://sigma-journal.com/2017/014/}\quad DOI: \url{https://doi.org/10.3842/SIGMA.2017.014}\par
\noindent\textbf{License and attribution:} Copyright retained by the named authors. Published by SIGMA under CC BY 4.0: \url{https://creativecommons.org/licenses/by/4.0/}. Source: \url{https://sigma-journal.com/about.html}. Adaptation consists of standalone excerpt formatting, cover and numbering restoration; original author language and mathematical text are retained.\par
\noindent\textbf{Editorial symbol notice (not author text):} Original author notation is preserved. The incidental D4 in Lemma4.9 character-group display denotes the explicitly defined D2 subgroup; the listed characters and restriction maps specify the calculation. Lemma5.6 expressly leaves other low-degree exact-sequence cases as routine calculations after a worked case.\par
\noindent\textbf{External original locator thm:main:} Original Theorem 1.1 summarizes the global wallpaper-action twist table; it is mentioned only in the general computation overview, outside the selected single p6m unit.\par
\medskip\hrule\medskip\noindent\textbf{Author text begins below.}\par
\setcounter{section}{1}
\setcounter{subsection}{0}
\setcounter{subsubsection}{0}
\setcounter{thm}{4}
\setcounter{equation}{0}
% Original source lines 330--348
To state our results in the `twisted' case, we introduce the following def\/inition for the point group $P$ of a $2$-dimensional space group $S$ that admits a non-trivial homomorphism $\phi\colon P \to \Z_2$.
\begin{itemize}\itemsep=0pt
\item
In the cases of \textsf{p2}, \textsf{p4} and \textsf{p6}, the point group $P$ is the cyclic group $\Z_{2m} = \langle C \,|\, C^{2m} \rangle$ of even order. We write $\phi_1\colon \Z_{2m} \to \Z_2$ for the unique non-trivial homomorphism given by $\phi_1(C) = -1$.

\item
In the other case, the point group $P$ is the dihedral group $D_n = \langle C, \sigma \,|\, C^n, \sigma^2, \sigma C \sigma C \rangle$
of degree $n$ and order $2n$, and $D_n$ is embedded into ${\rm O}(2)$ so that $C$ is a rotation of $\R^2$ and $\sigma$ is a ref\/lection. We def\/ine $\phi_0 \colon D_n \to \Z_2$ to be the composition of the inclusion $D_n \to {\rm O}(2)$ and $\det \colon {\rm O}(2) \to \Z_2$. Put dif\/ferently, $\phi_0(C) = 1$ and $\phi_0(\sigma) = -1$. This provides the unique non-trivial homomorphism $D_n \to \Z_2$ if $n$ is odd. In the case of even $n$, we def\/ine two more non-trivial homomorphisms $\phi_i \colon D_n \to \Z_2$ by
\begin{gather*}
\begin{cases}
\phi_1(C) = -1, \\
\phi_1(\sigma) = 1,
\end{cases} \qquad
\begin{cases}
\phi_2(C) = -1, \\
\phi_2(\sigma) = -1.
\end{cases}
\end{gather*}
\end{itemize}

\setcounter{section}{2}
\setcounter{subsection}{4}
\setcounter{subsubsection}{0}
\setcounter{thm}{4}
\setcounter{equation}{0}
% Original source lines 749--775
\section{The Leray--Serre spectral sequence and twists}\label{sec:spectral_seq_and_twist}

This section gives a geometric interpretation of the f\/iltration of $H^3_G(X; \Z)$ for the Leray--Serre spectral sequence through types of twists. This is carried out by identifying the Leray--Serre spectral sequence with another natural spectral sequence which computes the Borel equivariant cohomology.

Throughout this section, we assume that $G$ is a f\/inite group acting from the left on a `reasonable' space~$X$, such as a locally contractible, paracompact and regular topological space as in~\cite{FHT}, or a $G$-CW complex~\cite{May}.

\subsection{Spectral sequences}

The Borel equivariant cohomology $H^n_G(X; \Z)$ is def\/ined to be the (singular) cohomology of the quotient space $EG \times_G X$ of $EG \times X$ under the diagonal $G$-action $(\xi, x) \mapsto (\xi g, g^{-1}x)$, where $EG$ is the total space of the universal $G$-bundle $EG \to BG$. Associated to the f\/ibration $X \to EG \times_G X \to BG$ is the Leray--Serre spectral sequence
\begin{gather*}
E_r^{p, q} \Longrightarrow H^{p+q}_G(X; \Z)
\end{gather*}
converging to the graded quotient of a f\/iltration
\begin{gather*}
H^n_G(X; \Z) = F^0H^n_G(X; \Z) \supset F^1H^n_G(X; \Z) \supset \cdots \supset F^{n+1}H^n_G(X; \Z) = 0,
\end{gather*}
that is, $E_\infty^{p, q} = F^pH^{p+q}_G(X; \Z)/F^{p+1}H^{p+q}_G(X; \Z)$. The $E_2$-term is given by the group cohomology of~$G$
\begin{gather*}
E_2^{p, q} = H^p_{\mathrm{group}}(G; H^q(X; \Z)),
\end{gather*}
where the coef\/f\/icient $H^q(X; \Z)$ is regarded as a right $G$-module by the pull-back action. As a~convention of this paper, the group of $p$-cochains with values in a right $G$-module $M$ is denoted by $C^p_{\mathrm{group}}(G; M) = C(G^p, M) = \{ \tau \colon G^p \to M \}$, and the coboundary $\partial \colon C^p_{\mathrm{group}}(G; M) \to C^{p+1}_{\mathrm{group}}(G; M)$ is given by
\begin{gather*}
(\partial \tau)(g_1, \ldots, g_{p+1})= \tau(g_2, \ldots, g_{p+1})+ \sum_{i = 1}^p (-1)^i \tau(g_1, \ldots, g_ig_{i+1}, \ldots, g_{p+1}) \\
\hphantom{(\partial \tau)(g_1, \ldots, g_{p+1})=}{}+ (-1)^{p+1} \tau(g_1, \ldots, g_p)g_{p+1}.
\end{gather*}

As an application of the spectral sequence, we can obtain an identif\/ication $H^n_G(\pt; \Z) \cong H^n_{\mathrm{group}}(G; \Z)$. (We also have $H^n_{\mathrm{group}}(G; \Z) \cong H^{n-1}_{\mathrm{group}}(G; U(1))$ for $n \ge 2$ by the so-called exponential exact sequence.)

\setcounter{section}{4}
\setcounter{subsection}{0}
\setcounter{subsubsection}{0}
\setcounter{thm}{0}
\setcounter{equation}{0}
% Original source lines 1102--1211
\subsection{Some generality}\label{subsec:some_generality}

The cohomology $H^n(T^2; \Z)$ of the torus is well known, so that nothing remains to be proven in the case of \textsf{p1}.

For the point group $P$ of any $2$-dimensional space group, the vanishing $H^3(T^2; \Z) = 0$ implies $E^{0, 3}_\infty = 0$, so that
\begin{gather*}
H^3_P\big(T^2; \Z\big) = F^0H^3_P\big(T^2; \Z\big) = F^1H^3_P\big(T^2; \Z\big).
\end{gather*}
Note that each point group $P$ f\/ixes a point on $T^2$, so that
\begin{gather*}
F^3H^3_P\big(T^2; \Z\big) = H^3_P(\pt; \Z)= H^3_{\mathrm{group}}(P; \Z)= H^2_{\mathrm{group}}(P; U(1)).
\end{gather*}
Then the main task for the proof of Theorem \ref{thm:main} is to compute $H^3_P(T^2; \Z)$ and $F^2H^3_P(T^2; \Z)$, since in the case where $P$ is the cyclic group $\Z_n$ or the dihedral group $D_n$, the cohomology $H^m_P(\pt; \Z)$ is summarized as follows:
$$
\begin{array}{|c|c|c|c|c|}
\hline
P & H^0_P(\pt; \Z) & H^1_P(\pt; \Z) & H^2_P(\pt; \Z) & H^3_P(\pt; \Z) \tsep{2pt}\bsep{2pt}\\
\hline
\Z_n & \Z & 0 & \Z_n & 0 \\
\hline
D_n & \Z & 0 &
\begin{cases}
\Z_2 & (\mbox{$n$: odd}) \\
\Z_2^{\oplus 2} & (\mbox{$n$: even})
\end{cases} &
\begin{cases}
0 & (\mbox{$n$: odd}) \\
\Z_2 & (\mbox{$n$: even})
\end{cases}\tsep{10pt}\bsep{10pt}\\
\hline
\end{array}
$$
The degree $0$ part $H^0_P(\pt; \Z) = H^0(BP; \Z) = \Z$ is clear. Since $P$ is f\/inite, the degree $1$ part $H^1_P(\pt; \Z) \cong \operatorname{Hom}(P, \Z)$ is trivial. The degree $2$ part $H^2_P(\pt; \Z) \cong \operatorname{Hom}(P, U(1))$ can be seen by the classif\/ication of irreducible representations. Finally, the degree $3$ part $H^3_P(\pt; \Z) \cong H^2_{\mathrm{group}}(P; U(1))$ for $P = \Z_n, D_n$ can be found in \cite{Ka}.

In the rest of the section, we may use a structure of $T^2$ as a $P$-CW complex. In general, for a compact Lie group $G$, a \textit{$G$-CW complex} is an analogue of a CW complex made of \textit{$G$-cells}. A~$d$-dimensional $G$-cell is a $G$-space of the form $G/H \times e^d$, where $H \subset G$ is a closed subgroup and~$e^d$ is the standard $d$-dimensional cell. The $G$-action on $G/H$ is the left translation, whereas that on~$e^d$ is trivial. For the details, we refer the reader to \cite{May}.

We later compute a group cohomology via cohomology of a space:

\begin{lem} \label{lem:1_dim_complex}
Suppose that a finite group $G$ acts on a path connected space $Y$ fixing at least one point $\pt \in Y$. Suppose further that $Y$ is a CW complex consisting of only cells of dimension less than or equal to $1$. Then the following holds true for all $n \ge 0$.
\begin{gather*}
H^n_{\mathrm{group}}\big(G; H^1(Y; \Z)\big) \cong \tilde{H}^{n+1}_G(Y; \Z),
\end{gather*}
where $\tilde{H}^{n+1}_G(Y; \Z)$ stands for the reduced cohomology.
\end{lem}

Notice that a $G$-CW complex is naturally a CW complex.

\begin{proof}
Consider the Leray--Serre spectral sequence
\begin{gather*}
E_2^{p, q} = H^p_{\mathrm{group}}(G; H^q(Y; \Z))\Longrightarrow H^*_G(Y; \Z).
\end{gather*}
Note that $H^q(Y; \Z) = 0$ for $q \neq 0, 1$. The $E_2$-term $E_2^{p, 0} = H^p_{\mathrm{group}}(G; \Z)$ must survive into the direct summand $H^p_G(\pt; \Z)$ in $H^p_G(Y; \Z) \cong H^p_G(\pt; \Z) \oplus \tilde{H}^p_G(Y; \Z)$. Therefore it must hold that
\begin{gather*}
\tilde{H}^p_G(Y; \Z) \cong E^{p-1, 1}_\infty = E^{p-1, 1}_2\cong H^{p-1}_{\mathrm{group}}\big(G; H^1(Y; \Z)\big),
\end{gather*}
which completes the proof.
\end{proof}

We also prepare a simple lemma about group cohomology: Let $G$ be a f\/inite group, $c \colon G \to \Z_2 = \{ \pm 1 \}$ a surjective homomorphism, and $\tilde{\Z} = \Z_c$ the $G$-module such that its underlying group is $\Z$ and $G$ acts (from the right) by $m \mapsto m c(g)$. A typical example is a f\/inite subgroup $P \subset {\rm O}(2)$ such that $P \not\subset {\rm SO}(2)$ with $c$ the composition of the inclusion $P \to {\rm O}(2)$ and the determinant ${\rm O}(2) \to \Z_2$.

\begin{lem} \label{lem:group_coh_orientation_coeff}
Let $G$, $c$ and $\tilde{\Z}$ be as above. Then, $H^0_{\mathrm{group}}(G; \tilde{\Z}) = 0$ and $H^1_{\mathrm{group}}(G; \tilde{\Z}) \cong \Z_2$.
\end{lem}

\begin{proof}For any $n \in C^0_{\mathrm{group}}(G; \tilde{\Z}) = \Z$, its coboundary $\partial n\colon G \to \Z$ is $(\partial n)(g) = n(1 - c(g))$. Thus, the assumption that $c$ is surjective implies the vanishing of the $0$th cohomology. The inclusion $\operatorname{Ker}(c) \subset G$ induces an injection on $1$-cocycles
\begin{gather*}
Z^1_{\mathrm{group}}\big(G; \tilde{\Z}\big) \to Z^1_{\mathrm{group}}\big(\operatorname{Ker}(c); \tilde{\Z}\big) = \operatorname{Hom}(\operatorname{Ker}(c), \Z) = 0.
\end{gather*}
Thus, given a group $1$-cocycle $\phi \in Z^1_{\mathrm{group}}(G; \tilde{\Z})$, it holds that $\phi(g) = 0$ for all $g \in \operatorname{Ker}(c)$. If $g, h \not\in \operatorname{Ker}(c)$, then the cocycle condition $(\partial \phi)(g, h) = 0$ implies $\phi(g) = \phi(h)$. Therefore $\phi \colon G \to \Z$ is always of the form $\phi(g) = n (1 - c(g))/2$ for some $n \in \Z$. This provides the identif\/ication $Z^1_{\mathrm{group}}(G; \tilde{\Z}) \cong \Z$ as well as $B^1_{\mathrm{group}}(G; \tilde{\Z}) \cong 2\Z$, which completes the proof.
\end{proof}

In some cases, the computations of the Leray--Serre spectral sequence are similar, which we summarize as follows:

\begin{lem} \label{lem:typical_argument_non_orientation_preserving_case}
Let $G$ be a finite group acting on the torus $T^2$ such that:
\begin{itemize}\itemsep=0pt
\item there is a fixed point $\pt \in T^2$,

\item the $G$-action does not preserve the orientation of $T^2$.
\end{itemize}
Then the following holds true about the Leray--Serre spectral sequence:
\begin{itemize}\itemsep=0pt
\item[$(a)$] $F^2H^3_G(T^2; \Z) \cong E_2^{2, 1} \oplus E_2^{3, 0}$,

\item[$(b)$] $H^n_G(T^2; \Z) \cong \bigoplus_{p + q = n} E_2^{p, q}$ for $n \le 2$.
\end{itemize}
\end{lem}

\begin{proof}
In the Leray--Serre spectral sequence
$E_2^{p, q} = H^p_{\mathrm{group}}(G; H^q(T^2; \Z))$, the coef\/f\/icient in the group cohomology $H^0(T^2) \cong \Z$ is identif\/ied with the trivial $G$-module, and $H^2(T^2) \cong \Z$ with the $G$-module in Lemma \ref{lem:group_coh_orientation_coeff}. Then the relevant $E_2$-terms can be summarized as follows:
$$
\begin{array}{|c|c|c|c|c|c|}
\hline
q = 3 & 0 & 0 & 0 & 0 & 0 \\
\hline
q = 2 & 0 & \Z_2 & & & \\
\hline
q = 1 & E_2^{0, 1} & E_2^{1, 1} & E_2^{2, 1} & & \tsep{2pt}\bsep{2pt}\\
\hline
q = 0 & E_2^{0, 0} & E_2^{1, 0} & E_2^{2, 0} & E_2^{3, 0} & E_2^{4, 0} \tsep{2pt}\bsep{2pt}\\
\hline
E_2^{p, q} \tsep{2pt}\bsep{2pt} & p = 0 & p = 1 & p = 2 & p = 3 & p = 4\\
\hline
\end{array}
$$
Since $G$ f\/ixes $\pt \in T^2$, we have the decomposition $H^n_G(T^2; \Z) \cong H^n_G(\pt; \Z) \oplus \tilde{H}^n_G(T^2; \Z)$, where $\tilde{H}^n_G(T^2; \Z)$ is the reduced cohomology. Therefore the $E_2$-term $E_2^{n, 0} = H^n_{\mathrm{group}}(G; \Z) \cong H^n_G(\pt; \Z)$ must survive into the direct summand $H^n_G(\pt; \Z)$ in $H^n_G(T^2; \Z)$. This implies that $E_2^{n, 0} = E_\infty^{n, 0}$ is always a direct summand of the subgroups $F^pH^n_G(T^2; \Z) \subset H^n_G(T^2; \Z)$ and that $d_2 \colon E_2^{p-2, 1} \to E_2^{p, 0}$ is trivial. As a result, we get $E_2^{2, 1} = E_\infty^{2, 1}$ and the isomorphism (a). Also $E_2^{p, q} = E_\infty^{p, q}$ for $p + q \le 2$, and the isomorphism (b) follows.
\end{proof}

\setcounter{section}{4}
\setcounter{subsection}{4}
\setcounter{subsubsection}{0}
\setcounter{thm}{8}
\setcounter{equation}{0}
% Original source lines 1421--1873
\subsection{\textsf{p6m}}\label{subsec:p6m}

We let $\Pi = \Z a \oplus \Z b \subset \R^2$ be the lattice spanned by
$a = \left( \begin{matrix} 1 \\ 0 \end{matrix} \right)$ and
$b = \left( \begin{matrix} 1/2 \\ \sqrt{3}/2 \end{matrix} \right)$.
The point group $P$ is
$D_6 =\langle C, \sigma_1 \,|\, C^6, \sigma_1^2, \sigma_1C\sigma_1C \rangle=
\{1, C, C^2, C^3, C^4, C^5,\sigma_1, \sigma_2, \sigma_3, \sigma_4, \sigma_5, \sigma_6\}$, where $\sigma_\ell = C^{\ell - 1}\sigma_1$. This group acts on $\Pi$ and $\R^2$ through the inclusion $D_6 \subset {\rm O}(2)$ def\/ined by
\begin{gather*}
C = \left(
\begin{matrix}
1/2 & -\sqrt{3}/2 \\
\sqrt{3}/2 & 1/2
\end{matrix}
\right), \qquad \sigma_1 =
\left(
\begin{matrix}
1 & 0 \\
0 & -1
\end{matrix}
\right).
\end{gather*}
If we use the identif\/ications $a = 1$ and $b = \tau = \exp 2\pi i/6$ under $\R^2 = \C$, then the actions of $C \in D_6$ and $\sigma_1$ are given by the multiplication by $\tau$ and the complex conjugation, respectively. The closure of a fundamental domain is $\{ s a + t b \,|\, 0 \le s, t \le 1 \}$ or equivalently $\{ s + t \tau \,|\, 0 \le s, t \le 1 \}$. We decompose this region to def\/ine a $D_6$-CW decomposition of $T^2$ as follows:
$$
\begin{array}{|c|c|c|}
\hline
\mbox{0-cell} & \mbox{1-cell} & \mbox{2-cell} \\
\hline
\tilde{e}^0_0 = \pt &
\tilde{e}^1_{01} = (D_6/\{ 1, \sigma_1 \}) \times e^1 &
\tilde{e}^2 = D_6 \times e^2 \tsep{2pt}\\
\tilde{e}^0_1 = D^6/\{ 1, C^3, \sigma_1, \sigma_4 \} &
\tilde{e}^1_{02} = (D_6/\{ 1, \sigma_2 \}) \times e^1 & \\
\tilde{e}^0_2 = D^6/\{ 1, C^2, C^4, \sigma_2, \sigma_4, \sigma_6 \} &
\tilde{e}^1_{12} = (D_6/\{ 1, \sigma_4 \}) \times e^1 & \bsep{2pt}\\
\hline
\end{array}
$$

\begin{itemize}\itemsep=0pt
\item
($0$-cell)
The $0$-cell $\tilde{e}^0_0 = (D_6/D_6) \times e^0 = \pt$ is the unique f\/ixed point on $T^2$. The other $0$-cells are def\/ined as follows:
\begin{gather*}
\tilde{e}^0_1 =\left\{\frac{1}{2}, \frac{\tau}{2}, \frac{1+\tau}{2}\right\}\cong \big(D_6/\big\{ 1, C^3, \sigma_1, \sigma_4 \big\}\big) \times e^0, \\
\tilde{e}^0_2 =\left\{\frac{1 + \tau}{3}, \frac{2(1+\tau)}{3}\right\}\cong \big(D_6/\big\{ 1, C^2, C^4, \sigma_2, \sigma_4, \sigma_6 \big\}\big) \times e^0.
\\
\xygraph{
!{<0cm,0cm>;<1cm,0cm>:<0cm,1cm>::}
!{(0,0)}*+{\bullet}="a"
!{(2,0)}="b"
!{(1,1.7)}="c"
!{(3,1.7)}="d"
!{(1,0.6)}="e"
!{(2,1.2)}="f"
"a"-@{.}"b"
"a"-@{.}"c"
"b"-@{.}"c"
"b"-@{.}"d"
"c"-@{.}"d"
}
\qquad
\xygraph{
!{<0cm,0cm>;<1cm,0cm>:<0cm,1cm>::}
!{(0,0)}="a"
!{(2,0)}="b"
!{(1,1.7)}="c"
!{(3,1.7)}="d"
%!{(1,0.6)}="e"
%!{(2,1.2)}="f"
!{(1,0)}*+{\bullet}="g"
!{(0.5,0.85)}*+{\bullet}="h"
%!{(2,1.7)}="i"
%!{(2.5,0.85)}="j"
!{(1.5,0.85)}*+{\bullet}="k"
"a"-@{.}"g"
"g"-@{.}"b"
"a"-@{.}"h"
"h"-@{.}"c"
"b"-@{.}"k"
"k"-@{.}"c"
"b"-@{.}"d"
"c"-@{.}"d"
}
\qquad
\xygraph{
!{<0cm,0cm>;<1cm,0cm>:<0cm,1cm>::}
!{(0,0)}="a"
!{(2,0)}="b"
!{(1,1.7)}="c"
!{(3,1.7)}="d"
!{(1,0.6)}*+{\bullet}="e"
!{(2,1.2)}*+{\bullet}="f"
"a"-@{.}"b"
"a"-@{.}"c"
"b"-@{.}"c"
"b"-@{.}"d"
"c"-@{.}"d"
}
\end{gather*}

\item
($1$-cell)
For $0 \le i < j \le 2$, the $1$-cell $\tilde{e}^1_{ij}$ consists of the six segments connecting $\tilde{e}^0_i$ and $\tilde{e}^0_j$. They are of the forms $\tilde{e}^1_{01} = (D_6/\{ 1, \sigma_1 \}) \times e^1$, $\tilde{e}^1_{02} = (D_6/\{ 1, \sigma_2 \}) \times e^1$ and $\tilde{e}^1_{12} = (D_6/\{ 1, \sigma_4 \}) \times e^1$.
\begin{gather*}
\xygraph{
!{<0cm,0cm>;<1cm,0cm>:<0cm,1cm>::}
!{(0,0)}*+{\circ}="a"
!{(2,0)}*+{\circ}="b"
!{(1,1.7)}*+{\circ}="c"
!{(3,1.7)}="d"
%!{(1,0.6)}="e"
%!{(2,1.2)}="f"
!{(1,0)}*+{\circ}="g"
!{(0.5,0.85)}*+{\circ}="h"
%!{(2,1.7)}="i"
%!{(2.5,0.85)}="j"
!{(1.5,0.85)}*+{\circ}="k"
"a"-"g"
"g"-"b"
"a"-"h"
"h"-"c"
"b"-"k"
"k"-"c"
"b"-@{.}"d"
"c"-@{.}"d"
}
\qquad
\xygraph{
!{<0cm,0cm>;<1cm,0cm>:<0cm,1cm>::}
!{(0,0)}*+{\circ}="a"
!{(2,0)}*+{\circ}="b"
!{(1,1.7)}*+{\circ}="c"
!{(3,1.7)}*+{\circ}="d"
!{(1,0.6)}*+{\circ}="e"
!{(2,1.2)}*+{\circ}="f"
"a"-@{.}"b"
"a"-@{.}"c"
"b"-@{.}"c"
"b"-@{.}"d"
"c"-@{.}"d"
"a"-"e"
"b"-"e"
"c"-"e"
"d"-"f"
"b"-"f"
"c"-"f"
}
\qquad
\xygraph{
!{<0cm,0cm>;<1cm,0cm>:<0cm,1cm>::}
!{(0,0)}="a"
!{(2,0)}="b"
!{(1,1.7)}="c"
!{(3,1.7)}="d"
!{(1,0.6)}*{\circ}="e"
!{(2,1.2)}*{\circ}="f"
!{(1,0)}*+{\circ}="g"
!{(0.5,0.85)}*+{\circ}="h"
!{(2,1.7)}*+{\circ}="i"
!{(2.5,0.85)}*+{\circ}="j"
!{(1.5,0.85)}*+{\circ}="k"
"e"-"g"
"e"-"h"
"e"-"k"
"f"-"k"
"f"-"i"
"f"-"j"
"a"-@{.}"g"
"g"-@{.}"b"
"a"-@{.}"h"
"h"-@{.}"c"
"b"-@{.}"k"
"k"-@{.}"c"
"b"-@{.}"j"
"j"-@{.}"d"
"c"-@{.}"i"
"i"-@{.}"d"
}
\end{gather*}

\item
($2$-cell)
The $2$-cell $\tilde{e}^2 = D_6 \times e^2$ consists of the twelve small triangular regions surrounded by the $1$-cells.
\end{itemize}

Let $Y \subset T^2$ be the invariant subspace $Y = \tilde{e}^0_0 \cup \tilde{e}^0_1 \cup \tilde{e}^1_{01}$.


\begin{lem} \label{lem:p6m_Y}
The equivariant cohomology of $Y$ is given as follows:
$$
\begin{array}{|c|c|c|c|c|}
\hline
 & n = 0 & n = 1 & n = 2 & n = 3 \\
\hline
H^n_{D_6}(Y; \Z) & \Z & 0 & \Z_2^{\oplus 3} & \Z_2^{\oplus 2} \tsep{2pt}\bsep{2pt} \\
\hline
\end{array}
$$
\end{lem}

\begin{proof}
We can f\/ind $D_6$-invariant subspaces $U$ and $V$ in $Y$ which have the following equivariant homotopy equivalences
\begin{gather*}
U \simeq \tilde{e}^0_0 = \pt, \qquad
V \simeq \tilde{e}^0_1 = D_6/D_2, \qquad
U \cap V \simeq \tilde{e}^1_{01} \simeq D_6/\Z_2^{(1)},
\end{gather*}
where $D_2 = \{ 1, C^3, \sigma_1, \sigma_4 \}$ and $\Z_2^{(1)} = \{ 1, \sigma_1 \}$. The equivariant cohomology groups of these spaces can be summarized as follows:
$$
\begin{array}{|c|c|c|c|}
\hline
n = 3 & & \Z_2 \oplus \Z_2 & 0 \\
\hline
n = 2 & & \Z_2^{\oplus 2} \oplus \Z_2^{\oplus 2} & \Z_2 \tsep{2pt}\bsep{2pt}\\
\hline
n = 1 & & 0 & 0 \\
\hline
n = 0 & & \Z \oplus \Z & \Z \\
\hline
& H^n_{D_6}(Y) & H^n_{D_6}(U) \oplus H^n_{D_6}(V) & H^n_{D_6}(U \cap V)\tsep{2pt}\bsep{2pt}\\
\hline
\end{array}
$$
In the Mayer--Vietoris exact sequence
\begin{gather*}
\cdots \to H^n_{D_6}(Y) \to H^n_{D_6}(U) \oplus H^n_{D_6}(V) \overset{\Delta}{\to} H^n_{D_6}(U \cap V) \to H^{n+1}_{D_6}(Y) \to \cdots,
\end{gather*}
the homomorphism $\Delta$ is expressed as $\Delta (u, v) = j_U^*(u) - j_V^*(v)$ with $j_U \colon U \cap V \to U$ and $j_V \colon U \cap V \to V$ the inclusions. This immediately determines $H^0_{D_6}(Y) \cong \Z$ and $H^1_{D_6}(Y) = 0$. To complete the proof, we recall the identif\/ications
\begin{gather*}
H^2_{D_6}(U) \cong \operatorname{Hom}(D_6, U(1)) \cong \Z_2^{\oplus 2}, \qquad
H^2_{D_6}(V) \cong \operatorname{Hom}(D_2, U(1)) \cong \Z_2^{\oplus 2}, \\
H^2_{D_6}(U \cap V) \cong \operatorname{Hom}\big(\Z_2^{(1)}, U(1)\big) \cong \Z_2,
\end{gather*}
under which $j_U^*$ and $j_V^*$ are induced from the inclusions $D_2 \to D_6$ and $\Z_2^{(1)} \to D_6$. As a basis of~$H^2_{D_6}(U)$ we can choose the following $1$-dimensional representations $\rho_i\colon D_6 \to U(1)$ of $D_6$
\begin{gather*}
\rho_1 \colon \
\begin{cases}
C \mapsto 1, \\
\sigma_1 \mapsto -1,
\end{cases}
\qquad
\rho_2 \colon \
\begin{cases}
C \mapsto -1, \\
\sigma_1 \mapsto 1.
\end{cases}
\end{gather*}
Similarly, we can choose the following $1$-dimensional representations $\rho'_i$ of $D_2 = \{ 1, \sigma_1, C^3, \sigma_4 \}$ as a basis of $H^2_{D_6}(V)$
\begin{gather*}
\rho'_1 \colon \
\begin{cases}
C^3 \mapsto 1, \\
\sigma_1 \mapsto -1,
\end{cases}
\qquad
\rho'_2 \colon \
\begin{cases}
C^3 \mapsto -1, \\
\sigma_1 \mapsto 1.
\end{cases}
\end{gather*}
Now, we can see $H^2_{D_6}(Y) \cong \operatorname{Ker}\Delta \cong \Z_2^{\oplus 3}$, and it has the following basis
\begin{gather*}
\{ (\rho_1, \rho'_1), (\rho_2, \rho_2'), (0, \rho'_2) \} \subset \operatorname{Hom}(D_6, U(1)) \oplus \operatorname{Hom}(D_4, U(1)).
\end{gather*}
We can also see that $\Delta$ is surjective, and $H^3_{D_6}(Y) \cong \Z_2^{\oplus 2}$.
\end{proof}

Let $X_1$ be the $1$-skeleton of the $D_6$-CW complex $T^2$.

\begin{lem}
$H^3_{D_6}(X_1; \Z) \cong \Z_2^{\oplus 2}$.
\end{lem}

\begin{proof}
We cover $X_1 = \tilde{e}^0_0 \cup \tilde{e}^0_1 \cup \tilde{e}^0_2 \cup \tilde{e}^1_{01} \cup \tilde{e}^1_{02} \cup \tilde{e}^1_{12}$ by invariant subspaces $U'$ and $V'$ which admit the following equivariant homotopy equivalences
\begin{gather*}
U' \simeq Y, \qquad
V' \simeq \tilde{e}^0_2 = D_6/D_3, \qquad
U' \cap V' \simeq \tilde{e}^1_{02} \sqcup \tilde{e}^1_{12} \simeq D_6/\Z_2^{(2)} \sqcup D_6/\Z_2^{(4)},
\end{gather*}
where $D_3 = \{ 1, C^2, C^4, \sigma_2, \sigma_4, \sigma_6 \}$, $\Z_2^{(2)} = \{ 1, \sigma_2 \}$ and $\Z_2^{(4)} = \{ 1, \sigma_4 \}$. The equivariant cohomology groups of these spaces are summarized as follows:
$$
\begin{array}{|c|c|c|c|}
\hline
n = 3 & & \Z_2^{\oplus 2} \oplus 0 & 0\tsep{2pt}\bsep{2pt} \\
\hline
n = 2 & & \Z_2^{\oplus 3} \oplus \Z_2 & \Z_2 \oplus \Z_2 \tsep{2pt}\bsep{2pt}\\
\hline
n = 1 & & 0 & 0 \\
\hline
n = 0 & & \Z \oplus \Z & \Z \oplus \Z \\
\hline
& H^n_{D_6}(X_1) & H^n_{D_6}(U') \oplus H^n_{D_6}(V') & H^n_{D_6}(U' \cap V') \tsep{2pt}\bsep{2pt}\\
\hline
\end{array}
$$
The homomorphism $\Delta$ in the Mayer--Vietoris exact sequence
\begin{gather*}
\cdots \to H^2_{D_6}(X_1) \to H^2_{D_6}(U') \oplus H^2_{D_6}(V') \overset{\Delta}{\to} H^2_{D_6}(U' \cap V') \to H^3_{D_6}(X_1) \to \cdots
\end{gather*}
is expressed as $\Delta(u, v) = j_{U'}^*(u) - j_{V'}^*(v)$ by using the inclusions $j_{U'} \colon U' \cap V' \to U'$ and $j_{V'} \colon U' \cap V' \to V'$. An inspection proves that $j_{U'}^*$ agrees with the composition of the following two homomorphisms:
\begin{itemize}\itemsep=0pt
\item[(i)]
the inclusion that follows from the calculation of $H^2_{D_6}(Y)$ in Lemma \ref{lem:p6m_Y}
\begin{gather*}
H^2_{D_6}(U') \cong H^2_{D_6}(Y) \longrightarrow
\operatorname{Hom}(D_6, U(1)) \oplus \operatorname{Hom}(D_2, U(1)).
\end{gather*}

\item[(ii)]
the direct sum $i_2^* \oplus i_4^*$ of the homomorphisms
\begin{gather*}
i^*_2 \colon \ \operatorname{Hom}(D_6, U(1)) \to \operatorname{Hom}\big(\Z_2^{(2)}, U(1)\big), \qquad
i^*_4 \colon \ \operatorname{Hom}(D_2, U(1)) \to \operatorname{Hom}\big(\Z_2^{(4)}, U(1)\big),
\end{gather*}
induced from the inclusions $i_2\colon \Z_2^{(2)} \to D_6$ and $\Z_2^{(4)} \to D_2$.

\end{itemize}
Then, using the basis presented in the calculation of $H^2_{D_6}(Y)$, we f\/ind
\begin{gather*}
j_{U'}^*(\rho_1, \rho_1') = (\rho, \rho), \qquad j_{U'}^*(\rho_2, \rho_2') = (\rho, \rho), \qquad
j_{U'}^*(0, \rho_2') = (0, \rho),
\end{gather*}
where $\rho\colon \Z_2 \to \Z_2$ is the identity map generating $\operatorname{Hom}(\Z_2, U(1)) \cong \Z_2$. Hence $j_{U'}^*$ as well as $\Delta$ are surjective, and $H^3_{D_6}(X_1) \cong H^3_{D_6}(Y) \cong \Z_2^{\oplus 2}$.
\end{proof}

\begin{thm}[\textsf{p6m}]
$H^3_{D_6}(T^2; \Z) \cong \Z_2^{\oplus 2}$.
\end{thm}

\begin{proof}
The relevant part of the exact sequence for the pair $(T^2, X_1)$ is
\begin{gather*}
H^3_{D_6}\big(T^2, X_1\big) \to H^3_{D_6}\big(T^2\big) \to H^3_{D_6}(X_1) \to H^4_{D_6}\big(T^2, X_1\big).
\end{gather*}
By means of the excision axiom, we have $H^n_{D_6}(T^2, X_1) \cong H^{n-2}(\pt)$. Therefore we get $H^3_{D_6}(T^2) \cong H^3_{D_6}(X_1) \cong \Z_2^{\oplus 2}$.
\end{proof}

Let $\hat{\Z} = \Z_{\phi_1}$ be the $D_6$-module such that its underlying group is $\Z$ and $D_6$ acts via the homomorphism $\phi_1\colon D_6 \to \Z_2$ given by $\phi_1(C) = -1$ and $\phi_1(\sigma_1) = 1$.

\begin{lem} \label{lem:p6m_short_exact_sequence_of_modules}
There is an exact sequence of $D_6$-modules
\begin{gather*}
0 \to H^1\big(T^2; \Z\big) \to H^1(Y; \Z) \overset{\pi}{\to} \hat{\Z} \to 0
\end{gather*}
admitting a module homomorphism $s\colon \hat{\Z} \to H^1(Y; \Z)$ such that $\pi \circ s = 3$.
\end{lem}

\begin{proof}
Let $\eta_1, \eta_2 \in H_1(T^2)$ be the homology classes of the loops going along the vectors $1$ and $\tau$ respectively in the fundamental domain, which form a basis of $H_1(T^2) \cong \Z^2$. Also, let $\gamma_1, \gamma_2, \gamma_3 \in H_1(Y)$ be the homology classes of loops along $1$, $\tau$ and $\tau - 1$, which form a basis of $H_1(Y) \cong \Z^3$. The inclusion map $i\colon Y \to T^2$ relates these bases by $i_*\gamma_1 = \eta_1$, $i_*\gamma_2 = \eta_2$ and $i_*\gamma_3 = \eta_2 - \eta_1$. The actions of $C \in D_6$ and $\sigma_1$ on these bases are
\begin{gather*}
\begin{cases}
C_*\eta_1 = \eta_2, \\
C_*\eta_2 = \eta_2 - \eta_1,
\end{cases}
\qquad
\begin{cases}
{\sigma_1}_*\eta_1 = \eta_1, \\
{\sigma_1}_*\eta_2 = \eta_1 - \eta_2,
\end{cases}
\qquad
\begin{cases}
C_*\gamma_1 = \gamma_2, \\
C_*\gamma_2 = \gamma_3, \\
C_*\gamma_3 = - \gamma_1,
\end{cases}
\qquad
\begin{cases}
{\sigma_1}_*\gamma_1 = \gamma_1, \\
{\sigma_1}_*\gamma_2 = - \gamma_3, \\
{\sigma_1}_*\gamma_3 = - \gamma_2.
\end{cases}\!\!\!
\end{gather*}
Let $\{ h_1, h_2 \} \subset H^1(T^2)$ and $\{ g_1, g_2, g_3 \} \subset H^1(Y)$ be dual to the homology bases. They are related by $i^*h_1 = g_1 - g_3$ and $i^*h_2 = g_2 + g_3$, and the induced $D_6$-actions are as follows.
\begin{gather*}
\begin{cases}
C^*h_1 = - h_2, \\
C^*h_2 = h_1 + h_2,
\end{cases}\qquad
\begin{cases}
\sigma_1^*h_1 = h_1 + h_2, \\
\sigma_1^*h_2 = - h_2,
\end{cases}\qquad
\begin{cases}
C^*g_1 = - g_3, \\
C^*g_2 = g_1, \\
C^*g_3 = g_2,
\end{cases}\qquad
\begin{cases}
\sigma_1^*g_1 = g_1, \\
\sigma_1^*g_2 = - g_3, \\
\sigma_1^*g_3 = - g_2.
\end{cases}
\end{gather*}
These expressions allow us to prove that the cokernel of the homomorphism $i^* \colon H^1(T^2) \to H^2(Y)$ is isomorphic to $\hat{\Z}$, yielding the exact sequence. The homomorphism $s\colon \hat{\Z} \to H^1(Y)$ is given by $s(1) = g_1 - g_2 + g_3$.
\end{proof}

\begin{lem}
$H^n_{\mathrm{group}}(D_6; H^1(T^2; \Z)) = 0$ for $n = 0, 1, 2$.
\end{lem}

\begin{proof}
We use the long exact sequence in group cohomology induced from the exact sequence $0 \to H^1(T^2) \to H^1(Y) \overset{\pi}{\to} \hat{\Z} \to 0$ in coef\/f\/icients. By Lemmas~\ref{lem:1_dim_complex}, \ref{lem:p6m_Y} and~\ref{lem:cohomology_of_point_twisted_case} to be given in Section~\ref{sec:twisted_case}, we get the following:
$$
\begin{array}{|c|c|c|c|}
\hline
n = 2 & & \Z_2 & \Z_2 \\
\hline
n = 1 & & \Z_2 & \Z_2 \\
\hline
n = 0 & & 0 & 0 \\
\hline
& H^n_{\mathrm{group}}(D_6; H^1(T^2)) & H^n_{\mathrm{group}}(D_6; H^1(Y)) &
H^n_{\mathrm{group}}(D_6; \hat{\Z})\tsep{2pt}\bsep{2pt}\\
\hline
\end{array}
$$
It is clear that $H^0_{\mathrm{group}}(D_6; H^1(T^2)) = 0$. The homomorphism in group cohomology induced from $\pi\colon H^1(Y) \to \hat{\Z}$ is surjective in degree $1$ and $2$, because $\pi \circ s = 3$. This leads to the remaining vanishing.
\end{proof}

\begin{thm}[\textsf{p6m}]
The following holds true:
\begin{itemize}\itemsep=0pt
\item[$(a)$] $F^2H^3_{D_6}(T^2; \Z) \cong \Z_2$,

\item[$(b)$] the $D_6$-equivariant cohomology of $T^2$ in low degrees is as follows:
\begin{gather*}
H^0_{D_6}\big(T^2; \Z\big) \cong \Z, \qquad H^1_{D_6}\big(T^2; \Z\big) = 0, \qquad H^2_{D_6}\big(T^2; \Z\big) \cong \Z_2^{\oplus 2}.
\end{gather*}
\end{itemize}
\end{thm}

\begin{proof} In the $E_2$-term of the Leray--Serre spectral sequence $E_2^{p, q} = H^p_{\mathrm{group}}(D_6; H^q(T^2; \Z))$, the coef\/f\/icient $H^0(T^2)$ is identif\/ied with the trivial $D_6$-module $\Z$, and $H^2(T^2)$ with $\tilde{\Z}$ as in Lemma~\ref{lem:group_coh_orientation_coeff}. The group cohomology with coef\/f\/icients in $H^1(T^2)$ is already computed, and that in $\tilde{\Z}$ is also computed in Lemma~\ref{lem:group_coh_orientation_coeff}. The $E_2$-terms are summarized as follows:
$$
\begin{array}{|c|c|c|c|c|}
\hline
q = 3 & 0 & 0 & 0 & 0 \\
\hline
q = 2 & 0 & \Z_2 & & \\
\hline
q = 1 & 0 & 0 & 0 & \\
\hline
q = 0 & \Z & 0 & \Z_2^{\oplus 2} & \Z_2 \tsep{2pt}\bsep{2pt}\\
\hline
E_2^{p, q} & p = 0 & p = 1 & p = 2 & p = 3 \tsep{2pt}\bsep{2pt}\\
\hline
\end{array}
$$
This list and Lemma \ref{lem:typical_argument_non_orientation_preserving_case} lead to the theorem.
\end{proof}

\setcounter{section}{4}
\setcounter{subsection}{6}
\setcounter{subsubsection}{0}
\setcounter{thm}{15}
\setcounter{equation}{0}
% Original source lines 1921--1921
\section{The twisted case}\label{sec:twisted_case}

\setcounter{section}{5}
\setcounter{subsection}{2}
\setcounter{subsubsection}{0}
\setcounter{thm}{3}
\setcounter{equation}{0}
% Original source lines 2007--2086
\begin{lem} \label{lem:Z2_coefficient}
For any $\phi \colon G \to \Z_2$, there is a split exact sequence of groups
\begin{gather*}
0 \to H^n_G(X; \Z_\phi) \otimes \Z_2 \to H^n_G(X; \Z_2) \to \operatorname{Tor}\big(H^{n+1}_G(X; \Z_\phi), \Z_2\big) \to 0.
\end{gather*}
\end{lem}

\begin{proof}
For any homomorphism $\phi \colon G \to \Z_2$, let $(\Z_2)_\phi$ be the $G$-module such that its underlying group is $\Z_2$ and its $G$-action is given by $\phi \colon G \to \Z_2$. This $G$-module $(\Z_2)_\phi$ agrees with the trivial $G$-module $\Z_2$, even if $\phi$ is non-trivial. Then, looking at the cochain complexes def\/ining the equivariant cohomology, the usual proof of the universal coef\/f\/icient theorem leads to the lemma. Another proof is to use the Thom isomorphism, which unwinds the local coef\/f\/icients: Let $\underline{\R}_{\phi} \to X$ be the $G$-equivariant real line bundle on $X$ whose underlying bundle is $X \times \R$ and the action of $g \in G$ on $(x, r) \in X \times \R$ is $(x, r) \mapsto (gx, \phi(g)r)$. The Thom isomorphism theorem then provides
\begin{gather*}
H^n_G(X; A_\phi) \cong H^{n+1}_G(D, S; A),
\end{gather*}
where $A$ is $\Z_2$ or $\Z$, and $D \subset \underline{\R}_\phi$ and $S \subset \underline{\R}_\phi$ are the unit disk bundle and the unit sphere bundle, respectively. Then the usual universal coef\/f\/icient theorem leads to the present lemma.
\end{proof}

To compute the equivariant cohomology $H^n_G(X; \Z_\phi)$, we usually need the cohomology of the point $H^n_G(\pt; \Z_\phi)$. This cohomology is identif\/ied with the group cohomology $H^n_{\mathrm{group}}(G; \Z_\phi)$ by the degeneration of the Leray--Serre spectral sequence, but its direct computation is not realistic except for the simplest cases (cf.\ Lemma~\ref{lem:group_coh_orientation_coeff}). A useful way to compute it is:

\begin{lem} \label{lem:exact_seq}
For any $\phi \colon G \to \Z_2$, there are natural exact sequences
\begin{gather*}
\cdots \to H^{n-1}_G(\pt; \Z_\phi) \to H^n_G(\pt; \Z) \overset{i^*}{\to} H^n_{\operatorname{Ker}\phi}(\pt; \Z) \to H^n_G(\pt; \Z_\phi) \to \cdots, \\
\cdots \to H^{n-1}_G(\pt; \Z) \to H^n_G(\pt; \Z_\phi) \overset{i^*}{\to} H^n_{\operatorname{Ker}\phi}(\pt; \Z) \to H^n_G(\pt; \Z) \to \cdots,
\end{gather*}
where $i^*$ is induced from the inclusion $i \colon \operatorname{Ker}\phi \to G$.
\end{lem}

\begin{proof}
Let $G$ act on $\R_\phi = \R$ via $\phi \colon G \to \Z_2$. We can regard $\R_\phi$ as a $G$-equivariant real line bundle on $\pt$. We have the Thom isomorphisms
\begin{gather*}
H^n_G(\pt; \Z_\phi) \cong H^{n+1}_G(D, S; \Z), \qquad H^n_G(\pt; \Z) \cong H^{n+1}_G(D, S; \Z_\phi),
\end{gather*}
where $D$ and $S$ are the unit interval in $\R_\phi$ and its boundary, respectively. Note that $D$ is equivariantly contractible. Note also that $S \cong G/\operatorname{Ker}\phi$ as a $G$-space. Thus, we have isomorphisms
\begin{gather*}
H^n_G(D; \Z_\phi) \cong H^n_G(\pt; \Z_\phi), \qquad H^n_G(S; \Z_\phi) \cong H^n_{\operatorname{Ker}\phi}(\pt; \Z).
\end{gather*}
Substituting these isomorphisms and the Thom isomorphisms into the exact sequences for the pair $(D, S)$, we complete the proof.
\end{proof}

By means of the lemma above, we get:

\begin{lem} \label{lem:cohomology_of_point_twisted_case}
Let $P$ be $\Z_{2m}$, $D_{2m-1}$ or $D_{2m}$ with $m \ge 1$. The $P$-equivariant cohomology of the point with coefficients in $\Z_\phi$,
\begin{gather*}
H^n_P(\pt; \Z_\phi) \cong H^n_{\mathrm{group}}(P; \Z_\phi),
\end{gather*}
in low degrees is given as follows:
$$
\begin{array}{|c|c|c|c|c|c|}
\hline
P & \phi &
H^0_P(\pt; \Z_\phi) & H^1_P(\pt; \Z_\phi) &
H^2_P(\pt; \Z_\phi) & H^3_P(\pt; \Z_\phi) \tsep{2pt}\bsep{2pt}\\
\hline
\Z_{2m} & \phi_1 &
0 & \Z_2 & 0 & \Z_2 \\
\hline
D_{2m-1} & \phi_0 &
0 & \Z_2 & \Z_{2m-1} & \Z_2 \\
\hline
D_{2m} & \phi_0 &
0 & \Z_2 & \Z_{2m} & \Z_2^{\oplus 2} \tsep{2pt}\bsep{2pt}\\
\hline
D_{2m} & \phi_1, \phi_2 &
0 & \Z_2 & \Z_2 & \Z_2^{\oplus 2} \tsep{2pt}\bsep{2pt}\\
\hline
\end{array}
$$
\end{lem}

\begin{proof}
In the case of $D_{2m}$ with $m$ even and $\phi \neq \phi_0$, the f\/irst exact sequence in Lemma~\ref{lem:exact_seq} leads to $H^0_{D_{2m}}(\pt; \Z_\phi) = 0$ and $H^1_{D_{2m}}(\pt; \Z_\phi) \cong \Z_2$. This computation also shows that $H^2_{D_{2m}}(\pt; \Z_\phi)$ contains $\Z_2$ as a subgroup. Here, applying the universal coef\/f\/icient theorem to $H^n_P(\pt; \Z)$, we compute the cohomology with coef\/f\/icients in $\Z_2$ to have $H^1_{D_{2m}}(\pt; \Z_2) \cong \Z_2^{\oplus 2}$. If we apply the universal coef\/f\/icient theorem in Lemma \ref{lem:Z2_coefficient} to $H^n_P(\pt; \Z_\phi)$, then
\begin{gather*}
H^1_{D_{2m}}(\pt; \Z_2) \cong \Z_2 \oplus \operatorname{Tor}\big(H^2_{D_{2m}}(\pt; \Z_\phi), \Z_2\big).
\end{gather*}
Thus the consistency of these computations implies $H^2_{D_{2m}}(\pt; \Z_\phi) \cong \Z_2$. Based on this result, the second sequence in Lemma~\ref{lem:exact_seq} suggests that $H^3_{D_{2m}}(\pt; \Z_\phi)$ is either $\Z_2^{\oplus 2}$ or $\Z_4$. If we apply the universal coef\/f\/icient theorem to $H^n_P(\pt; \Z)$, then $H^2_{D_{2m}}(\pt; \Z_2) \cong \Z_2^{\oplus 3}$. If we compute this cohomology applying Lemma~\ref{lem:Z2_coefficient} to $H^n_P(\pt; \Z_\phi)$, then
\begin{gather*}
H^2_{D_{2m}}(\pt; \Z_2) \cong \Z_2 \oplus \operatorname{Tor}\big(H^3_{D_{2m}}(\pt; \Z_\phi), \Z_2\big).
\end{gather*}
Therefore we conclude that $H^3_{D_{2m}}(\pt; \Z_\phi), \Z_2) \cong \Z_2^{\oplus 2}$ by the consistency. In the other cases, a~combined use of the two exact sequences in Lemma~\ref{lem:exact_seq} determines the group $H^n_P(\pt; \Z_\phi)$ without dif\/f\/iculty.
\end{proof}
\begin{thebibliography}{99}
\bibitem[7]{FHT}
Freed D.S., Hopkins M.J., Teleman C., Loop groups and twisted {$K$}-theory~{I},
 \href{https://doi.org/10.1112/jtopol/jtr019}{\textit{J.~Topol.}} \textbf{4} (2011), 737--798, \href{http://arxiv.org/abs/0711.1906}{arXiv:0711.1906}.

\bibitem[14]{Ka}
Karpilovsky G., Projective representations of f\/inite groups, \textit{Monographs
 and Textbooks in Pure and Applied Mathematics}, Vol.~94, Marcel Dekker, Inc.,
 New York, 1985.

\bibitem[19]{May}
May J.P., Cole M., Comezana G.R., Costenoble S.R., Elmendorf A.D., Greenlees
 J.P., Lewis L.G., Piacenza R.J., Triantaf\/illou G., Waner S., Equivariant
 homotopy and cohomology theory, \href{https://doi.org/10.1090/cbms/091}{\textit{CBMS Regional Conference Series in
 Mathematics}}, Vol.~91, Amer. Math. Soc., Providence, RI, 1996.

\end{thebibliography}
\end{document}
